%\ssscp{Etyma distributed across multiple subgroups}{15-04-03-01}
In several chapters (\ref{ch:CM}--\ref{ch:Ban})
we occasionally referred to PCM (Proto
\ili{Central} Maluku) etyma, noting that these were
geographic terms-of-convenience with widely
varying distributions, and that there was no proto-language based on the historical
phonology to which they could be assigned.
In \srf{15-05} onwards we consider how these cognates came to
be distributed across multiple subgroups.
A selection of etyma widely distributed in central
Maluku are given in \trf{15-03a} and \trf{15-03b}.

\begin{table}[h]
	\centering\tcp{Etyma widely distributed in central Maluku (Part 1)}{15-03a}
	{\wg\st{2pt}\begin{tabularx}{\linewidth}
	{lll>{\raggedright\arraybackslash}X}
	\lsptoprule
Etymon	&	Gloss	&	Subgroups	&	Notes	\\	\midrule
**batəlu	&	thick 	&	SB, AS, 	&	\hp{\,}	\\	\hhline{~}
&	(dimension)	& STB, \tsc{Cel} &	\hp{\,}	\\	\hhline{~}
&		&		&	\hp{\,}	\\	\hhline{~}
&		&		&	\multirow{-4}{=}{\raggedright probably ultimately linked to PMP *təbəl through patterns of metathesis widespread in the region. cf. \tsc{POc} *matolu thick	}	\\
**dada(n)	&	thousand	&	SB, AS &	replacement of PMP *Ribu	\\	\hhline{~}
&		&	(\ili{Kayeli})		&	\hp{\,}	\\
**gəgəl	&	armpit	&	\tsc{Tal}, SB, 	&	cf. \tsc{PAnD} *gəlgəl notch	\\	\hhline{~}
&		&	AS, STB	&	\hp{\,}	\\
**gəlan	&	tree sp, 					&	SB, AS	&	\hp{\,}	\\	\hhline{~}
				&	\it{Melaleuca}		&					&	\multirow{-2}{=}{\raggedright source of commercially important tea tree oil}	\\	\hhline{~}
				&	\it{leucodendron}	&					&	\hp{\,}	\\
**labu	&	ashes, dust	&	SB, AS, 	&	sporadic alongside PMP *qabu	\\	\hhline{~}
&		&	STB, TB	&	\hp{\,}	\\
**lambu	&	shirt	&	SB, AS	&	cf. \tsc{PPh} *lambuŋ robe, long garment	\\
**maluku	&	eel	&	SB, AS, STB	&	eels figure prominantly in origin stories	\\
**maRan	&	mother's 	&	SB, AS	&	\hp{\,}	\\	\hhline{~}
&	natal clan	&&	\hp{\,}	\\
**maRuka	&	girl, maiden	&	\mc{2}{l}{SB, AS,	Banda}	\\
**ma-siRi	&	drunk, dizzy	&	SB, AS	&	\hp{\,}	\\
**mbətu	&	night, 24-hour 	&	\tsc{Tal}, SB, AS	&	\hp{\,}	\\	\hhline{~}
&day cycle&&	\multirow{-2}{=}{\raggedright Also possibly STB languages \ili{Kei} and Yamdena}	\\
**muyu	&	\tsc{neg}, no, not	&	SB, AS 	&	post-predicate in all languages	\\	\hhline{~}
&		&	(\ili{Alune})	&	\hp{\,}	\\
%**nakaw	&	steal	&	\tsc{Tal}, SB, AS, STB (\tsc{E. Seram})	&	ultimately from PMP *takaw	\\
\lspbottomrule
\end{tabularx}}
\end{table}

\begin{table}[h]
	\centering\tcp{Etyma widely distributed in central Maluku (Part 2)}{15-03b}
	{\wg\st{2pt}\begin{tabularx}{\linewidth}
	{lll>{\raggedright\arraybackslash}X}
	\lsptoprule
Etymon	&	Gloss	&	Subgroups	&	Notes	\\	\midrule
**nawa	&	arenga 	&	SB, AS	&	ultimately from PMP *qanahaw	\\	\hhline{~}
&	palm tree	&		&	\hp{\,}	\\
**ndau	&	brother-	&	SB, AS	&	often sister's husband	\\	\hhline{~}
&	in-law	&		&	\hp{\,}	\\
**niwəR	&	coconut	&	SB, AS, \ili{Banda}	&	irregular from PMP *niuR	\\
**ntibu	&	fly (v.)	&	SB, AS, 	&	STB only \tsc{E. Seram} and \tsc{E. Islands}	\\	\hhline{~}
&		&	\ili{Banda}, STB	&	\hp{\,}	\\
**ŋaRək	&	year	&	SB, AS, STB	&	possibly PMP *taqun shifted to `full'	\\
**ŋiRut	&	mucous	&	\tsc{Tal}, SB, 	&	\hp{\,}	\\	\hhline{~}
&		&	AS, \ili{Banda}	&	\hp{\,}	\\
**sabe	&	buy	&	SB, AS	&	\hp{\,}	\\
**sa(R)i	&	paddle (verb)	&	SB, AS	&	cf. PMP *bəRsay	\\
**səgət	&	high tide	&	\tsc{Tal}, SB, AS	&	\hp{\,}	\\
**sibaw	&	cuscus pouch, 	&	SB, AS	&	\hp{\,}	\\	\hhline{~}
				&	female cuscus	&&	\hp{\,}	\\
**sombal	&	sail (v.), &	SB, AS, STB, 	&	loan from \ili{Makasar} \it{sombalaʔ}	\\	\hhline{~}
&	go by boat		&TB, \tsc{SSul}&	\hp{\,}	\\
**ta(q)un	&	full	&	SB, AS	&	reflexes of *ma-pənuq `full' demoted to minor role, but often co-exist	\\
**tanəi	&	work	&	SB, AS	&	\hp{\,}	\\
**təwa	&	know, 	&	SB, AS	&	\hp{\,}	\\	\hhline{~}
&	understand	&&	\hp{\,}	\\
**tubuRa, &	spit, spray &	\tsc{Tal}, SB, AS, 	&	\hp{\,}	\\	\hhline{~}
**batuRa	&medicinal herbs	&\ili{Banda}, STB&	\multirow{-2}{=}{\raggedright ultimately from PMP *buRah; both forms interspersed}	\\
**utum	&	hundred	&	SB, AS, \ili{Banda}	&	replaces PMP *Ratus	\\
\lspbottomrule
\end{tabularx}}
\end{table}

The nodes listed in Tables \ref{tab:15-03a} and \ref{tab:15-03b}
reflect the limitations of our awareness.
It is not unlikely that reflexes are also found in other nodes
of which we are not aware.\footnote{
		We would like to thank a reviewer who pointed the following additional
		distributions that were not in the original submission: **sombal
		‘sail’ (\trf{15-03b}) also occurs in South Sulawesi,
		STB, and TB; **batəlu ‘thick’ (\trf{15-03a}) also occurs in Celebic;
		**pəsu ‘fart’ (\trf{15-04}) also occurs in STB; **seru ‘weave’
		(\trf{15-05}) also occurs in the \tsc{\ili{Bungku}-Tolaki} node of \tsc{Celebic};
		and **ləmba ‘carry on pole’ (\trf{15-07}) also occurs in \tsc{South Sulawesi} and \ili{Sumbawa}.
		The \tsc{South Sulawesi} cognates of **sombal
		and **ləmba were first noted by \citet[68, 70]{Mills1981}.}
In addition,
these terms may not necessarily be found in all branches of the subgroups listed.
Occasionally we give a PCMP or PCEMP reconstruction made by \citet{BlustTrussel2020}
in the notes column,
to indicate that others have also identified the widespread nature
of certain cognate sets.
This does not indicate that we accept PCMP or PCEMP as valid nodes.

There are terms reconstructed for Proto-\tsc{Rote-Meto} (PRM)
also found in \ili{Bima} (\tsc{East Sumbawa} node of \tsc{\ili{Bima}-Lembata})
and often other nodes of \tsc{\ili{Bima}-Lembata}.
\citet{VerheijenGrimes1995} note that after 1727 \ili{Bima} exerted
its influence into western Flores and beyond.
Selected examples are shown in \trf{15-04}.
Some of these also have cognates in \tsc{\ili{Central} Timor} and/or
nodes of \tsc{Timor-Babar} in addition to \tsc{Rote-Meto}.
The notes column lists the nodes in which cognates are currently known to occur.
See the entries in \citet{Edwards2021} for full details.
(\tsc{Flores} covers both \tsc{Flores-Lembata} (FL) with \tsc{W.
Flores} and/or \tsc{C. Flores}).

\begin{table}[h]
	\centering\tcp{Selected Proto-Rote-\ili{Meto} etyma also found in Bima}{15-04}
	\wg\st{2pt}\begin{tabularx}{\linewidth}{ll>{\raggedright\arraybackslash}X}\lsptoprule
PRM	&	gloss	&	notes\\	\midrule
**ɓee	&	where; in, at	&	\ili{Bima} \it{ɓee} ‘where’\\
**deŋe	&	kapok tree, Bombax ceiba	&	\ili{Bima} \it{riŋi} ‘kapok’; \ili{Helong} \it{deŋe}\\
**əndi	&	bring, take	&	\ili{Bima}, \tsc{Flores}, \tsc{Sumba-Havu}; \ili{Helong}, \tsc{E. Timor}, \ili{Galolen}; CT; \ili{Sumbawa} [smw] \it{ənti}\\
**hao	&	feed	&	\ili{Bima} \it{pao} ‘mouthful, place in mouth’\\
**hesu	&	fart	&	\ili{Bima}, \tsc{Sumba-Havu}; \ili{Helong},  \tsc{E. Timor}; STB; PCMP **pəsu\\
**kera	&	brother-in-law	&	\ili{Bima}, \tsc{Flores}, \tsc{Sumba}; \ili{Helong}; CT\\
**mafo	&	cool, shadow	&	\ili{Bima}, \tsc{Sumba-Havu}; \tsc{E. Timor}\\
**mofu	&	fall	&	\ili{Bima} \it{mou} ‘fall’; S. \ili{Mambae} (CT) \it{mou}\\
**ndoro	&	walk around	&	\ili{Bima}, \ili{Havu}; \ili{Helong}\\
**ŋgasi	&	shout, speak	&	\ili{Bima}, WC. \tsc{Flores}; \tsc{SW. Maluku}; CT\\
**sarakaen	&	sand	&	\ili{Bima}, \tsc{\ili{Havu}-Dhao}; \tsc{\ili{Lakalei}-Idate}; CT; \ili{Kamang} (non-\tsc{An}) \it{siraiŋ}\\
**selut	&	replace, exchange	&	\ili{Bima}, \tsc{Sumba}; \ili{Helong}, \tsc{E. Timor} (potential cognates in Java, Bali, Sulawesi)\\
**taɗu	&	opposite	&	\ili{Bima} \it{ntando} ‘facing one another’\\
**toka	&	impede, against	&	\ili{Bima}, \tsc{Flores-Lembata}; \ili{Helong}\\
**tumbi	&	k.o. tree with soft wood	&	\ili{Bima} \it{tuwi} ‘kind of soft wood’\\
		\lspbottomrule
	\end{tabularx}
\end{table}

Terms reconstructed for Proto-\tsc{Rote-Meto} with cognates in
languages of Flores (but not \ili{Bima} or Sumba) are shown in \trf{15-05}.
Travelling north-south between Timor
and Flores requires deliberate effort
and is at odds with the prevailing seasonal \ili{east}-west winds.

\begin{table}[p]\centering\tcp{Selected Proto-Rote-\ili{Meto} etyma also found in Flores}{15-05}
\begin{adjustbox}{totalheight=\textheight}
	\st{2pt}\wg\begin{tabularx}{1.07\linewidth}{ll>{\raggedright\arraybackslash}X}\lsptoprule
PRM			&	gloss	&	notes\\	\midrule
**biti	&	slingshot	&	FL; \ili{Helong}\\
**ɓoŋgo	&	gourd, pumpkin	&	\tsc{C. Flores}; \ili{Helong}, \tsc{E. Timor}\\
**ɓusa	&	dog	&	FL; \tsc{E. Timor}; CT\\
**deki	&	kiss	&	FL; \ili{Helong}, \tsc{E. Timor}\\
**ɗəma	&	deep; drown	&	\tsc{Flores}, \ili{Havu} (cf. STB **l[əa]ma ‘inside’, see \trf{11-02})\\
**ɗoi		&	carry on shoulder 	&	\tsc{C. Flores}, \tsc{\ili{Havu}-Dhao}\\	\hhline{~}
&	with pole	&	\hp{\,}	\\
**feku	&	flute	&	\tsc{C. Flores}\\
**imun	&	mosquito, midge	&	\tsc{C. Flores}\\
**kteom	&	sea urchin	&	\tsc{W. Flores}; \ili{Helong}, \tsc{E. Timor} (cognates in Philippines, Pre-RM **(ka)-tayum)\\
**kepe	&	tick	&	\tsc{C. Flores}\\
**koɗe	&	monkey	&	\tsc{WC. Flores}\\
**kodo	&	bird (generic)	&	\tsc{Flores}, \ili{Havu}; \ili{Oirata} (non-\tsc{An}) \it{olo} ‘quail’\\
**kura	&	scorpion	&	FL, \ili{Havu}; \ili{Helong}, \tsc{E. Timor}, \ili{Galolen}, \tsc{SW. Maluku} (pre-RM **sakuraŋ)\\
**loɗo	&	straight	&	FL, \ili{Havu}; \ili{Helong}, \tsc{SW. Maluku}\\
**lolo	&	beam, stick	&	FL; \ili{Helong}, \tsc{SW. Maluku}\\
**lutu	&	fine, tiny, small	&	FL, \ili{Dhao}; \ili{Helong}, \tsc{E. Timor}\\
**malus	&	betel pepper	&	FL; \tsc{E. Timor}, \tsc{Wetar}, \tsc{SW. Maluku}; CT\\
**mbasa	&	slap	&	FL; \ili{Helong}, \tsc{E. Timor}, \tsc{Wetar}, \tsc{SW. Maluku}; CT\\
**muse	&	seed, pip	&	\tsc{C. Flores}, \ili{Havu}; \tsc{E. Timor}\\
**nafuʔ	&	body hair, feather	&	FL; \tsc{E. Timor} (pre-\ili{Meto} **rabu(k))\\
**nitas	&	\it{Sterculia foetidea}	&	\tsc{Flores}; \tsc{E. Timor}\\	\hhline{~}
&	(tree sp.)	&	\hp{\,}	\\
**rutus	&	rust	&	FL, \ili{Havu}; \tsc{E. Timor}\\
**seru	&	weave; sword	&	\hp{\,}	\\	\hhline{~}
&	 of loom	&	\hp{\,}	\\	\hhline{~}
&			&	\multirow{-3}{=}{\raggedright FL, \ili{Havu}; \ili{Helong}, \tsc{E. Timor} (cf. Proto-\ili{Bungku}- \ili{Tolaki} *horu \citep[437]{Mead1998} and Proto-\ili{Ambon} **sedu \citep[18, 93, 100]{Stresemann1927}}	\\
**simbo	&	receive, accept	&	\tsc{Flores}, \tsc{\ili{Havu}-Dhao}; \tsc{E. Timor}; CT; SB\\
**soɗa	&	sing	&	\tsc{C. Flores}, \ili{Havu}\\
**soi		&	open, release;	&	\hp{\,} \\	\hhline{~}
&	 redeem, ransom	&	\multirow{-2}{=}{\raggedright \tsc{Flores}; \ili{Helong}, \tsc{E. Timor} (often replaces PMP *təbus ‘ransom, redeem’)}	\\
**tamae	&	bedbug	&	\tsc{Flores}; \ili{Helong} (pre-RM **tamayuŋ)\\
**tema	&	whole, entire	&	FL, \ili{Havu}; \ili{Helong}, \tsc{E. Timor}, \tsc{SW. Maluku}\\
**tetu	&	upright; midday	&	\tsc{Flores}, \ili{Havu}; \ili{Helong}, \tsc{E. Timor}\\
**tuku	&	row	&	FL, \ili{Havu}; \ili{Helong} (cf. \tsc{POc} *tokon ‘punting pole, pole a boat in shallow water’)\\
		\lspbottomrule
	\end{tabularx}
\end{adjustbox}
\end{table}

Terms reconstructed for Proto-\tsc{Rote-Meto} with cognates in
Sumba are shown in \trf{15-06}.
In recent times, \ili{Dhao} -- a member of \tsc{Sumba-Havu} -- has regular contact with Rote.
Sumba and \ili{Havu} do not.
But since both \ili{Havu}
and Rote share lontar-palm economies which figure prominently
in both cultures \citep{Fox1977}, contact in the past could have been quite different.
%Selected examples are shown in \trf{15-06}.

\begin{table}[p]
	\centering\tcp{Selected Proto-Rote-\ili{Meto} etyma also found in \tsc{Sumba}}{15-06}
	\wg\st{4pt}\begin{tabularx}{\linewidth}{ll>{\raggedright\arraybackslash}X}
\lsptoprule
PRM			&	gloss												&	notes\\	\midrule
**afa		&	insipid, bland, 						&	\tsc{Sumba}\\	\hhline{~}
				&	tasteless										&	\hp{\,}	\\
**ɓanafi&	sea cucumber								&	\tsc{Sumba}, FL; \ili{Helong}, \tsc{E. Timor} (commercially important, possibly borrowed into some Australian languages, \citealp[445]{Schapper2022}) \\
**ɓoŋgo	&	round												&	\tsc{Sumba}; \ili{Helong}, \tsc{E. Timor}\\
**ɗeta	&	dip (e.g. tip	of finger)		&	\tsc{Sumba-Havu}\\
**ɗolu	&	draw water,  								&	\tsc{Sumba-Havu}\\	\hhline{~}
				&	lower; fish (v.)						&	\hp{\,}	\\
**ɗoʔi	&	lever, pry out							&	\tsc{Sumba}, FL; \ili{Helong}\\
**fa		&	no; just										&	\tsc{Sumba-Havu}\\
**fake	&	\it{Parkia speciosa}, 			&	\tsc{Sumba}; \ili{Helong}, \tsc{E. Timor}\\	\hhline{~}
				&	bitter bean									&	\hp{\,}	\\
**fandu	&	dry season									&	\tsc{Sumba-Havu}\\
**ʔfenu	&	\it{Aleurites moluccanus}, 	&	\tsc{Sumba}, \tsc{WC. Flores} (pre-\ili{Meto} **ka-felu)\\	\hhline{~}
				&	candlenut										&	\hp{\,}	\\
**haŋga	&	hand span, span							&	\tsc{Sumba-Havu}, FL; (pre-RM **paŋga, cf. PMP *zaŋkal)\\
**lifu	&	pond, billabong, 						&	\hp{\,}	\\	\hhline{~}
				&	pool, lake									&	\hp{\,}	\\	\hhline{~}
				&															&	\multirow{-3}{=}{\raggedright \tsc{Sumba}, FL; \ili{Helong}, \tsc{E. Timor}; CT (cf. \ili{Buru} \it{libu} ‘sinkhole’, PCMP **tibu ‘pool, pond, lagoon’)}	\\
**ŋgola	&	k.o. lizard									&	\tsc{Sumba-Havu}, \tsc{C. Flores}\\
**oka		&	pen, corral									&	\tsc{Sumba-Havu}; \ili{Helong}\\
**poke	&	blind												&	\tsc{Sumba-Havu}; \tsc{Wetar}; STB (only \ili{Sekar} \it{pok{\tl}pok} ‘blind’)\\
**rose	&	rub, wipe										&	\tsc{Sumba-Havu}, \tsc{Flores}\\
**seɗa	&	k.o. red coral 							&	\tsc{Sumba-Havu}\\	\hhline{~}
				&	used as a bead							&	\hp{\,}	\\
**tana	&	cover												&	\tsc{Sumba-Havu}; \ili{Helong} (pre-RM **taŋa)\\
**tanaʔ	&	thorn												&	\tsc{Sumba}, \tsc{E. Timor}, \tsc{SW. Maluku}; CT (pre-\ili{Meto} **tarak)\\
**tuku	&	throw												&	\tsc{Sumba}\\
		\lspbottomrule
	\end{tabularx}
\end{table}

There are other terms which occur in \tsc{Timor-Babar},
as well as more distant subgroups,
such as \tsc{Sula-Buru}, \tsc{\ili{Ambon}-Seram}, and/or \tsc{Oceanic}.
A selection of such words are given \trf{15-07}.
These words also have widely varying distributions,
and appear to not be found in some Wallacean subgroups.
The nodes listed in \trf{15-07} reflect the limitations of our awareness.
Where \citet{BlustTrussel2020} have made a PCMP
and PCEMP reconstruction, this is given in the notes column of \trf{15-07}.
Where \citet{Edwards2021} has made a Proto-\tsc{Rote-Meto} reconstruction
this is also given.
The reader can consult \citet{Edwards2021} and/or \citet{BlustTrussel2020}
for more details on the reflexes their distribution.

\begin{table}[p]
	\wg\centering\tcp{Other cognate sets distributed unevenly across Wallacea}{15-07}
	\begin{adjustbox}{totalheight=\textheight}
	\st{2pt}\begin{tabularx}{1.05\linewidth}{llll>{\raggedright\arraybackslash}X}\lsptoprule
Etymon		&	PRM				&	gloss								&	nodes												&	notes\\	\midrule
**buqi		&	**fui			&	pour, douse					&	BL, TB, STB, \tsc{Oc}				&	PCEMP *buqi\\
**buu			&	**fuu			&	blow								&	BL, TB, \tsc{Tal}, \tsc{Oc}	&	Blust PCMP *buu; onomatopoeic\\
**kambeR	&	**kambe		&	saliva							&	TB, CT, AS									&	\hp{\,}	\\
**koti		&	**koti		&	cut									&	TB, \tsc{Oc}								&	PCMP *koti\\
**kumu(l)	&	**kumu		&	wild pigeon					&	TB, SB, \tsc{Oc}						&	\tsc{POc} *g(a)umu\\
**kupu		&	**kupu		&	fog									&	TB, CT, \tsc{Oc}						&	\tsc{POc} *kupu(k)\\
**lemba		&	**lemba		&	carry on pole				&	BL, TB, CT,  								&	\hp{\,}	\\	\hhline{~}
					&						&											&	SB, AS, \tsc{SSul}, 				&	\multirow{-2}{=}{\raggedright	likely loan from \ili{Makasar} \it{lembaraʔ}}	\\	\hhline{~}
					&						&											&	\ili{Sumbawa} 										&	\hp{\,}	\\
**lemur		&	**lemuk		&	dolphin							&	TB, \tsc{Aru}								&	\hp{\,}	\\
**ləsi		&						&	excess							&	TB, SB, AS									&	PCMP *lesi\\
**lətay		&	**lete		&	above, 							&	BL, TB, AS, STB							&	PCMP *letay\sub{1}\\	\hhline{~}
					&						&	mountain						&															&	\hp{\,}	\\
**lutuR		&	**lutu		&	stone barrier				&	BL, TB, CT, 								&	see \citet{Schapper2020a}\\	\hhline{~}
					&						&											&	AS, STB, \tsc{Aru}					&	\hp{\,}	\\
**malip		&	**malis		&	laugh, smile				&	BL, TB, \tsc{Aru}, 					&	PCEMP **malip; \hp{\,}	\\	\hhline{~}
					&						&											&	\tsc{Tal}, SB, AS, STB, 		&	see \frf{15-09}\\	\hhline{~}
					&						&											&	\mc{2}{l}{\ili{Banda}, SHWNG, \tsc{Oc}}	\\
**ma-putu	&	**mbutu		&	hot									&	BL, TB, SB, AS							&	\hp{\,}	\\
**maRus		&	**maus		&	tame, docile				&	BL, TB, SB, AS, \tsc{Oc}		&	\hp{\,}	\\
**maya		&	**maa			&	tongue							&	BL, TB, \tsc{Tal}, SB, 			&	PCEMP **maya;	\\	\hhline{~}
					&						&											&	AS, STB, \ili{Banda}, 						&	see \frf{15-08}	\\	\hhline{~}
					&						&											&	\mc{2}{l}{SHWNG, \tsc{Oc}}	\\
**(m)butu	&	**futu, 	&	group, bunch 				&	BL, TB, SB, AS,							&	\hp{\,}	\\	\hhline{~}
					&	**mbutu		&											&	STB, \tsc{Oc}								&	\multirow{-2}{=}{\raggedright Maluku `ten', Flores `four'; PCEMP **butu}\\
**ŋəli		&	**ŋoli		&	canine tooth, 			&	BL, TB, AS									&	\hp{\,}	\\	\hhline{~}
					&						&	tusk								&															&	\hp{\,}	\\
**sa(q)uR	&	**soo			&	sew									&	BL, TB, CT, 								&	PCMP *sora\\	\hhline{~}
					&						&											&	STB, SB, AS									&	\hp{\,}	\\
**salili	&	**salili	&	armpit							&	BL, TB, AS, 								&	\hp{\,}	\\	\hhline{~}
					&						&											&	\ili{Sumbawa}											&	\hp{\,}	\\
**taduku	&	**taruku	&	chiton  						&	TB, \tsc{Oc}								&	\tsc{P-E. Oc} *tadruku\\
%					&						&	\it{Polyplacophora}	&															&	\hp{\,}	\\	\hhline{~}
**talada	&	**talaɗa	&	middle, centre			&	BL, TB, AS									&	\hp{\,}	\\
**tuir		&	**tuin		&	follow							&	BL, TB, AS									&	\hp{\,}	\\
**usi			&	**usi			&	CCC, PPP						&	BL, TB, SB, \tsc{Tal}				&	PCMP **usi	\\
**wasi		&	**osi			&	garden							&	TB, CT, SB, AS							&	\hp{\,}	\\
**zuzu		&	**ɗuɗu		&	tinder, oakum				&	BL, TB, \tsc{SSul}					&	\hp{\,}	\\
		\lspbottomrule	\end{tabularx}
	\end{adjustbox}
\end{table}
